%----------------------------------------------------------------------
(Full group-theoretic analysis: \S\ref{app:banach-tarski}.)
\section{Set-Theoretic Consequences}\label{sec:settheory}
%----------------------------------------------------------------------

We now construct standard set-theoretic structure from the\footnote{Section-level Toulmin. \emph{Move}: derive set-theoretic axioms as consequences of the framework. \emph{ZFC comparison}: ZFC adopts separation, foundation, infinity, and choice as independent axioms targeted at specific pathologies. Here, Russell exclusion follows from profile linearity (adopted for algebraic faithfulness), foundation from profile linearity + finite $\kappa$, and infinity/choice from the operationalist axiom (adopted for engineering soundness). The framework provides structural \emph{diagnoses}: self-reference produces an affine offset; infinite chains exceed finite $\kappa$; actual/potential infinity are separated by no discriminator at any level of the halting tower. \emph{Consequential elegance}: three independently motivated principles handle what ZFC needs four separate axioms for.}
discrimination framework (algebra + regulatory axioms),
noting where axioms of ZFC appear as consequences and
where the framework diverges.

\subsection{Extensionality}

\begin{definition}[$\kappa$-equivalence]
Two collections $A$ and $B$ are \emph{$\kappa$-equivalent}, written
$A \equiv_\kappa B$, if for every discriminator $d$ in $\kappa$ and
every element $a$ in the ambient population,
$(\tau_d(a), \sigma_d(a))$ computed over $A$ equals
$(\tau_d(a), \sigma_d(a))$ computed over $B$.
\end{definition}

\begin{theorem}[Extensionality]\label{thm:ext}\footnote{Axiom class \textsf{A} $\cdot$ depth 1.  $\kappa$-equivalence is defined by discriminator profiles, and extensionality is immediate from the definition. Finite-$\kappa$: yes. Re-description (standard, in discrimination language).}
$\kappa$-equivalence is the identity criterion for sets:
$A \equiv_\kappa B$ if and only if no discriminator in $\kappa$
distinguishes $A$ from $B$. This is not an axiom but a direct
consequence of the definition of $\kappa$-equivalence as the
identity relation on equivalence classes.
\end{theorem}

\begin{proof}
Sets \emph{are} $\kappa$-equivalence classes (by construction).
Two classes are equal if and only if they agree on every discriminator.
This is the definition of equivalence-class identity.
\end{proof}

\begin{remark}
Extensionality in ZFC states that two sets are equal iff they have the
same members. Here the ``same members'' condition is relativized to
$\kappa$: two sets are equal \emph{under a given discrimination regime}.
A finer $\kappa$ may distinguish sets that a coarser $\kappa$ identifies.
Extensionality is $\kappa$-relative, which is strictly more informative
than the ZFC version.
\end{remark}

\subsection{Comprehension and Russell's predicate}

\begin{definition}[Discriminative comprehension]\label{def:comprehension}
For any predicate $P$ that can be articulated as a discriminator,
the act of articulation is a $\kappa$-evolution producing a new
discriminator $d_P$. The set $\{a \in S : P(a)\}$ is the collection
$\{a \in S : \tau_{d_P}(a) = 1\}$ (or $\sigma_{d_P}(a) = 1$,
depending on the polarity convention).

Comprehension is unrestricted in the sense that any articulable
predicate induces a $\kappa$-evolution. It is restricted by the
algebra: if the resulting wedge product is degenerate, the
comprehension yields the trivial set or the empty residue.
\end{definition}

\begin{theorem}[Exclusion of Russell's predicate]\label{thm:russell}\footnote{Axiom class \textsf{AEP} $\cdot$ depth 3 $\cdot$ \S\ref{app:russell-detail}.  \emph{Warrant}: profile linearity requires linear maps; $\tau_{{d_P}}(0) = 1 \neq 0$ is affine. \emph{Backing}: Def~\ref{def:profile-linear}, Def~\ref{def:eval-linear}. \emph{Qualifier}: this is a relocation of restriction, not an elimination (Remark 4.8). \emph{Rebuttal}: ZFC handles this via the separation axiom; here it follows from profile linearity, at the cost of also requiring evaluation linearity. Scope of exclusion: 8 of 256 functions admissible for $n=3$. Finite-$\kappa$: yes. Novel (bilinear form mechanism). Model: \S\ref{sec:concrete-model}. (\S4.6). [O1]}
The predicate $P(x) = (x \notin x)$ does not produce a non-trivial
set under discriminative comprehension.
\end{theorem}

\begin{proof}
We construct an explicit bilinear form whose diagonal
gives self-membership and whose antisymmetrization is the
wedge product.

\emph{Step 1: The profile pairing.}
By profile linearity (Definition \ref{def:profile-linear}),
the maps $\tau, \sigma: V \to \GF(2)^S$ (sending each
discriminator to its $\tau$- or $\sigma$-profile on $S$)
are linear. When the population $S$ contains elements
that are also discriminators (i.e.\ $S \cap V \neq
\varnothing$), we can define the \emph{profile pairing}:
\begin{equation}\label{eq:profile-pairing}
  B: V \times V \to \GF(2), \qquad
  B(x, y) = \tau_x(y),
\end{equation}
the $\tau$-profile of discriminator $x$ evaluated at
element $y$ (where $y$ is also in $V$). By profile
linearity, $B$ is linear in the first argument:
$B(x_1 + x_2, y) = \tau_{x_1 + x_2}(y) = \tau_{x_1}(y)
+ \tau_{x_2}(y) = B(x_1, y) + B(x_2, y)$.

For $B$ to be bilinear, it must also be linear in the
second argument: $B(x, y_1 + y_2) = B(x, y_1) + B(x, y_2)$.
This requires two conditions: (a) the population $S$
inherits the $\GF(2)$-vector space structure of $V$ when
$S \cap V \neq \varnothing$ (so that $y_1 + y_2$ is
well-defined as a population element), and (b) each
$\tau_x: S \to \GF(2)$ acts linearly on this inherited
structure. Condition (a) holds by construction when
discriminators are elements of $S$. Condition (b) is an
extension of profile linearity (Definition \ref{def:profile-linear}) to the evaluation argument:
we require not only that the map $x \mapsto \tau_x$ is
linear, but that each $\tau_x$ is itself a linear
functional on $V$. We adopt this as an axiom:

\begin{definition}[Evaluation linearity]\label{def:eval-linear}
For each discriminator $x \in V$, the $\tau$-profile
$\tau_x: V \to \GF(2)$ is a $\GF(2)$-linear functional
when restricted to $S \cap V$. Together with profile
linearity (Definition \ref{def:profile-linear}), this
makes the profile pairing $B(x,y) = \tau_x(y)$ a
bilinear form on $V$.

This is a second regulatory axiom, not a free consequence
of the framework. Its motivation: if discriminator
profiles were non-linear in their evaluation argument,
then composing discriminators (the wedge product) would
not faithfully represent composed discriminations ---
the algebraic structure of $\bigwedge V$ would be
decoupled from the operational behavior of the
discrimination process. Evaluation linearity (Definition \ref{def:eval-linear}) is the
axiom that keeps the algebra honest.

\emph{What evaluation linearity excludes.}
Any predicate whose $\tau$-profile is a non-linear
function of the evaluation argument is inadmissible.
Concretely:

\emph{Excluded}: (i) The Russell predicate
($\tau_{d_P}(x) = 1 + x(x)$, affine). (ii) Threshold
predicates: ``$a$ has property $\tau$ iff at least two
of three conditions hold'' --- this is a degree-2
polynomial over $\GF(2)$, not linear. (iii) Any
predicate defined by a Boolean function of degree
$\geq 2$ in the evaluation argument.

\emph{Preserved}: (i) All single-condition predicates:
``$a$ has property $\tau$ iff condition $C$ holds'' where
$C$ is a single discriminator evaluation. These are
linear by construction. (ii) XOR-type predicates:
``$a$ has $\tau$ iff an odd number of conditions hold''
--- the sum over $\GF(2)$ of individual conditions is
linear. (iii) All predicates obtained by composition of
linear predicates via $\wedge$ and $+$.

Over $\GF(2)$, the linear functions $V \to \GF(2)$ are
exactly the elements of $V^*$ (the dual space). With
$\dim(V) = n$, there are $2^n$ linear functionals and
$2^{2^n}$ total functions $V \to \GF(2)$. Evaluation
linearity restricts to the linear ones: a fraction
$2^n / 2^{2^n}$ of all possible predicates. For $n = 3$:
8 of 256 functions are admissible (3.1\%). The exclusion
is severe but systematic: it excludes all non-linear
Boolean functions while preserving the full dual space.
\end{definition}

\emph{Step 2: Self-membership is the diagonal.}
Self-membership of $x$ (the question ``is $x \in x$?'')
is $\tau_x(x) = B(x, x)$: the profile pairing evaluated
on the diagonal.

\emph{Step 3: The antisymmetrization is the wedge product.}
The antisymmetrization of $B$ is
$\widetilde{B}(x, y) = B(x, y) - B(y, x) = B(x, y)
+ B(y, x)$ (since $-1 = 1$ in $\GF(2)$). This is
$\tau_x(y) + \tau_y(x)$: the XOR of $x$'s $\tau$-verdict
on $y$ and $y$'s $\tau$-verdict on $x$. This
antisymmetrized form is the wedge product's evaluation:
$\widetilde{B}(x, y)$ vanishes when $x = y$ (since
$\tau_x(x) + \tau_x(x) = 0$ in $\GF(2)$), which is the
defining property $x \wedge x = 0$.

\emph{Step 4: The Russell predicate is affine, not linear.}
The Russell predicate $P(x) = (x \notin x)$ requires a
discriminator $d_P$ whose $\tau$-profile is
$\tau_{d_P}(x) = 1 + \tau_x(x) = 1 + B(x,x)$
(the negation in $\GF(2)$ of self-membership).

Over $\GF(2)$, the Frobenius endomorphism $x \mapsto x^2$
is the identity ($x^2 = x$ for all $x \in \GF(2)$), so
$B(x,x)$ is actually $\GF(2)$-linear in $x$ (not
quadratic in the classical sense): $B(x{+}y, x{+}y) =
B(x,x) + B(x,y) + B(y,x) + B(y,y)$, and the cross-terms
$B(x,y) + B(y,x) = \widetilde{B}(x,y)$ vanish on any
pair where the antisymmetrization vanishes. The diagonal
$Q(x) = B(x,x)$ is therefore a linear functional on $V$.

However, $\tau_{d_P}(x) = 1 + Q(x)$ is \emph{affine},
not linear: it has a constant offset of $1$. Profile
linearity (Definition \ref{def:profile-linear}) requires
that $\tau_d$ be a linear map $V \to \GF(2)$, and linear
maps must satisfy $\tau_d(0) = 0$. But
$\tau_{d_P}(0) = 1 + B(0,0) = 1 + 0 = 1 \neq 0$. The
Russell predicate's profile sends the zero discriminator
to $1$, which no linear functional does.

Therefore $d_P$ violates profile linearity (Definition \ref{def:profile-linear}) and is not a
valid grade-1 discriminator. The Russell predicate is
articulable (it triggers a $\kappa$-evolution), but the
resulting profile is affine, not linear. No grade-1
discriminator with a linear $\tau$-profile can carve out
the Russell set. The predicate is algebraically
ill-formed --- not prohibited by fiat, but excluded by the
linearity structure that the framework requires of its
discriminators.
\end{proof}

\begin{remark}[Profile linearity as the load-bearing axiom]
The dissolution depends on profile linearity
(Definition \ref{def:profile-linear}), not on
$x \wedge x = 0$ directly. The wedge product's role is
indirect: $x \wedge x = 0$ ensures that $B(x,x)$ is
linear over $\GF(2)$, but the Russell predicate's
constant offset $1 + B(x,x)$ is affine. The dissolution
is: self-reference produces an affine map; discriminators
require linear maps; affine $\neq$ linear.
(Full bilinear form construction: \S\ref{app:russell-detail}.)
\end{remark}

\begin{remark}[The restriction is relocated]
Profile linearity (Definition \ref{def:profile-linear}) is a \emph{regulatory axiom}: it
restricts what counts as a valid discriminator, just as
ZFC's separation axiom restricts what counts as a valid
set-formation. The Russell predicate is not
``algebraically degenerate'' in the exterior algebra ---
it is \emph{excluded by the linearity requirement on
$\tau$-profiles}. This is a lateral move from ZFC, not a
dissolution: the restriction has been relocated from set
formation to discriminator admissibility.

The relocation is arguably more natural --- profile
linearity is a single structural axiom that handles
Russell, foundation, and the bilinear form structure
simultaneously, whereas ZFC requires separate axioms for
each pathology. But it is still a restriction. The
framework does not conjure paradox-freedom from pure
algebra; it achieves it by requiring discriminators to be
linear. The paper should be read with this understanding:
where a restriction appears, it has been
internalized into the discriminator definition,
and the paradox reduces to a violation of that definition.
The restriction is relocated, not removed.
\end{remark}

\begin{remark}[General pathology protection]
One may ask whether the framework might be vulnerable to
set-theoretic pathologies beyond Russell's. The framework
handles pathologies via two mechanisms:
\begin{itemize}[nosep]
  \item \textbf{Profile linearity}
    (Definition \ref{def:profile-linear}): restricts
    discriminator admissibility. This excludes the Russell
    predicate (affine, not linear) and bounds comprehension.
  \item \textbf{The exterior algebra}: $x \wedge x = 0$
    ensures self-composition is degenerate, and the
    four-valued logic is paraconsistent ($p \wedge \neg p$
    produces $(1,1)$, a residue, not explosion).
\end{itemize}
Specific pathologies are handled by: comprehension
bounded by profile linearity (Definition \ref{def:profile-linear}), foundation by
linear independence of the profile matrix
(Theorem \ref{thm:foundation}), and choice scoped to
Boolean dimensions with the infinite case depending on
the operationalist axiom (Definition \ref{def:operationalist}) (Theorem \ref{thm:choice}).
These are structural restrictions, not ad hoc patches,
but they \emph{are} restrictions. The framework does not
achieve paradox-freedom from pure algebra alone; it
achieves it from the algebra together with profile
linearity and the operationalist axiom.
\end{remark}

\subsection{Boolean structure as a local achievement}

\begin{definition}[Boolean dimension]\label{def:boolean-dim}
A dimension $d$ of $\kappa$ is \emph{Boolean over a collection $S$}
if for every $a \in S$, $\chi_d(a) = 1$ --- that is, exactly one of
the two predicates fires for every element ($\tau_d(a) = 1$
and $\sigma_d(a) = 0$, or vice versa). The two independent
partitions happen to be complementary on $S$: every element
$\tau$-comprehended is $\sigma$-uncomprehended, and vice versa.
No gaps (both silent) and no overlaps (both firing). The law of
excluded middle holds locally on dimension $d$.
\end{definition}

\begin{proposition}[Non-default Booleanness]\label{prop:nondefault}
Booleanness of a dimension $d$ over $S$ is not guaranteed by the
framework. It is an additional structural condition that must hold
for the specific discriminator and collection. The default state
--- independent $\tau$ and $\sigma$ with possible gaps and
overlaps --- is strictly weaker than Boolean partition.
\end{proposition}

\begin{proof}
The membership profile takes values in $\GF(2) \times \GF(2)$,
a space with four elements. Boolean partition requires confinement
to the subset $\{(1,0), (0,1)\}$, which is a proper subset.
Nothing in the definition of discriminator (Definition
\ref{def:discriminator}) enforces this confinement.
\end{proof}

\begin{remark}
This is perhaps the central philosophical claim: classical logic's
excluded middle is not a law of thought but a \emph{theorem}, local
to specific $\kappa$-dimensions over specific collections, and it
must be verified constructively. Boolean complementation requires a
mechanism; it is not free.
\end{remark}

\subsection{Foundation and well-foundedness}

\begin{theorem}[Partial foundation]\label{thm:foundation}\footnote{Axiom class \textsf{AP} $\cdot$ depth 4 $\cdot$ \S\ref{app:foundation-full}.  \emph{Warrant}: self-membership excluded by profile linearity; chain length bounded by finiteness of $\kappa$. \emph{Qualifier}: weaker than ZFC foundation (permits cycles $\geq 2$). \emph{Rebuttal}: ZFC excludes all cycles via a separate axiom; here self-loops follow from profile linearity, at the cost of permitting longer cycles (Remark~\ref{rem:cycles}). $\in$-induction: resolved. Cycles consume $\leq k{-}1$ dimensions; induct on $\dim(V/\kappa)$ (Appendix \ref{app:o2-cycles}). [O2 resolved.] Finite-$\kappa$: yes. Novel (intermediate position: no self-loops, permits cycles). Model: \S\ref{sec:concrete-model}.}
\begin{enumerate}[label=(\roman*),nosep]
  \item Self-membership ($a \in a$) is excluded by
    profile linearity (Definition \ref{def:profile-linear}) (Theorem \ref{thm:russell}).
  \item For any finite $\kappa$-state, only finitely many
    membership relations can be witnessed. The depth of the
    membership relation is bounded by $\dim(\kappa)$.
\end{enumerate}
\end{theorem}

\begin{proof}
(i) is Theorem \ref{thm:russell}: self-membership
requires an affine profile, which violates profile
linearity.

(ii) Each membership link $a \in_{d,\kappa} B$ is
witnessed by a discriminator $d \in \kappa$. With
$\dim(\kappa) = n$, there are $n$ basis discriminators
and $2^n - 1$ non-zero elements of $\bigwedge^1 V$. Each
can witness at most finitely many membership relations
(bounded by $|S|^2$). Therefore the total number of
membership links that $\kappa$ can witness is finite.
(Full analysis: \S\ref{app:foundation-full}.)
\end{proof}

\begin{remark}[Cycles of length $\geq 2$ are permitted]
\label{rem:cycles}
The framework does \emph{not} exclude membership cycles
of length $\geq 2$. A cycle $a \in b \in a$ can be
witnessed by two independent discriminators $d_1, d_2$
with $\tau_{d_1}(a) = 1$, $\sigma_{d_1}(b) = 1$,
$\tau_{d_2}(b) = 1$, $\sigma_{d_2}(a) = 1$. This is
algebraically consistent: $d_1 \wedge d_2 \neq 0$ (since
$d_1, d_2$ are independent), and no profile linearity (Definition \ref{def:profile-linear})
violation occurs. The two predicates fire independently
on $a$ and $b$, with no self-reference.

This is a genuine structural difference from ZFC, where
the axiom of regularity prevents all cycles. In the
discrimination framework, only self-membership (length 1)
is excluded. Longer cycles are permitted because they
do not involve self-reference --- each link is witnessed
by a separate discriminator acting on distinct elements.

The algebraic obstruction to self-membership is
$x \wedge x = 0$ (the diagonal of the profile pairing
vanishes), which prevents an element from bearing a
discriminative relation to itself. But
$d_1 \wedge d_2 \neq 0$ for independent $d_1, d_2$, so
two-element cycles face no such obstruction.
\end{remark}

\begin{remark}
In ZFC, Foundation and the axiom of dependent choice
together prevent both infinite descending chains and
finite cycles. The framework prevents infinite chains
(by finiteness of $\kappa$) and self-membership (by
profile linearity (Definition \ref{def:profile-linear})), but permits finite cycles. This more
permissive structure may be natural for applications
where circular reference is meaningful (e.g.\ mutual
recursion, cyclic data structures, or self-referential
systems that reference each other without referencing
themselves).

Within the landscape of non-well-founded set theories,
the framework's position is: strictly weaker than ZFC's
foundation (which excludes all cycles), and also distinct
from Aczel's Anti-Foundation Axiom (AFA), which permits
all non-well-founded sets via a universal coalgebra.
The framework excludes self-loops (length 1) but permits
all longer cycles, without the systematic treatment of
non-well-founded sets that AFA provides. This is an
intermediate position that has not, to our knowledge,
been previously studied.

We trace the consequences of permitting cycles through
the paper's main structures:

\emph{Extensionality.} Cycles do not threaten
extensionality. If $a \in b \in a$, then $a$ and $b$
have distinct $\kappa$-profiles (witnessed by independent
discriminators $d_1, d_2$). They are distinguishable,
hence not $\kappa$-equivalent. The membership graph may
have cycles, but the equivalence relation on objects
does not.

\emph{Homology.} A membership cycle $a \in b \in a$
is a cycle in the \emph{membership graph}, not in the
\emph{chain complex}. The corresponding 2-chain
$d_1 \wedge d_2$ has boundary
$\partial(d_1 \wedge d_2) = d_1 + d_2 \neq 0$
(since $d_1, d_2$ are independent). The membership
cycle is not a homological cycle: it does not
represent a gap in the discrimination complex.
The two notions of ``cycle'' are distinct.

\emph{Sheaf conditions.} Covering the regime by
sub-regimes containing $d_1$ and $d_2$ respectively:
the sheaf condition requires compatibility on overlaps.
Since $d_1$ and $d_2$ are independent, their overlap
is the trivial sub-regime, and the sheaf condition is
automatically satisfied. Membership cycles do not
produce sheaf obstructions.

\emph{$\in$-induction.} This is where cycles have
genuine consequences. Standard $\in$-induction
(``if the property holds for all members of $x$, it
holds for $x$'') fails in the presence of cycles: for
$a \in b \in a$, establishing the property for $a$
requires it for $b$, which requires it for $a$. The
resulting proof obligation is circular, not
well-founded. A weaker form --- $\in$-induction modulo
the cycle equivalence relation (identify elements in
the same cycle) --- may still hold, but this has not
been verified. Whether useful forms of $\in$-induction
survive in the framework is an open question.
\end{remark}

\begin{remark}
In ZFC, foundation is an axiom imposed to prevent pathologies
such as $x \in x$. Here, $x \in x$ is already excluded by
$x \wedge x = 0$ (Theorem \ref{thm:russell}), and the more
general well-foundedness condition follows from the irreversibility
of $\kappa$-evolution. Foundation is a consequence of profile linearity (Definition \ref{def:profile-linear})
(excluding self-membership) and finiteness of $\kappa$
(bounding chain length). It is not an independent
assumption in the sense that it follows from axioms
already adopted for other purposes; but it does depend
on profile linearity, which is itself a regulatory axiom.
\end{remark}

\subsection{Empty set and pairing}

\begin{proposition}[Empty set]\label{prop:empty}
For any discriminator $d$ and collection $S$, the sub-collection
$\{a \in S : \tau_d(a) = 1 \text{ and } \sigma_d(a) = 1
\text{ and } \tau_d(a) = 0\}$ (a contradictory filter) yields
the empty collection. More naturally, a $\kappa$-deformation
that projects away all dimensions yields the trivially empty result.
The empty set exists as a possible outcome of filtration.
\end{proposition}

\begin{proposition}[Pairing]\label{prop:pairing}
Given elements $a, b$ in a collection $S$, there exists a
$\kappa$-evolution that introduces a discriminator $d$ for which
$\tau_d(a) = 1$, $\tau_d(b) = 1$, and $\tau_d(c) = 0$ for all
$c \in S \setminus \{a,b\}$. The set $\{a,b\}$ is the $\tau$-class
of $d$. This follows from the ability to articulate arbitrary
predicates: $a$ and $b$ form a $\kappa$-equivalence class:
choose $\kappa$ with $a + b \in \kappa^\perp$,
so every $d \in \kappa$ satisfies $d(a) = d(b)$,
while other elements are separated
(Appendix \ref{app:pairing-linearity}),
hence it triggers a $\kappa$-evolution.
\end{proposition}

\subsection{Union and the role of $\GF(2)$}

\begin{proposition}[Symmetric difference as primitive]\label{prop:symdiff}
Given two sets $A$ and $B$ defined as $\tau$-classes of
discriminators $d_A$ and $d_B$ respectively, the natural
algebraic combination $\tau_{d_A}(a) + \tau_{d_B}(a)$ in
$\GF(2)$ yields the \emph{symmetric difference}
$A \triangle B$, not the union $A \cup B$. This is because
$1 + 1 = 0$ in $\GF(2)$: elements in both $A$ and $B$ are
annihilated.
\end{proposition}

\begin{corollary}
Boolean union $A \cup B = (A \triangle B) \triangle (A \cap B)$
is a derived operation requiring explicit construction of the
intersection $A \cap B$. In the discrimination framework,
intersection (the overlap cell) is logically prior to union.
\end{corollary}

\begin{remark}
This reversal of the usual ZFC presentation --- where union is
axiomatic and intersection is defined via comprehension --- is
a direct consequence of working over $\GF(2)$ rather than
Boolean algebra.
\end{remark}

\subsection{Power set}

\begin{definition}[Discriminative power set]\label{def:powerset}
The \emph{power set} of a collection $S$ under regime $\kappa$ is
the set of all sub-collections of $S$ obtainable by
$\kappa$-deformations (projections to sub-regimes). If $\kappa$
has $n$ dimensions, the power set has at most $2^n$ elements.
\end{definition}

\begin{remark}
The power set is not a fixed Platonic totality. It is bounded by
the current $\kappa$-state and grows as $\kappa$ evolves. This
gives power sets a constructive character absent from ZFC, where
$\mathcal{P}(\mathbb{N})$ has a fixed (if undetermined)
cardinality.
\end{remark}

\begin{conjecture}[Constructive continuum]\label{conj:CH}
In the discrimination framework, the ``cardinality'' of
the power set of $\mathbb{N}$ is not a fixed value but an
open-ended function of the discrimination regime, growing
with each new predicate articulated over $\mathbb{N}$.
This suggests a correspondence with the independence of
the continuum hypothesis from ZFC. However, the
independence of CH is proved via forcing models (Cohen)
and inner models (G\"{o}del), and establishing a genuine
correspondence would require showing that no
$\kappa$-evolution can witness either CH or $\neg$CH, and
that this non-witnessability maps onto the specific model-
theoretic independence results. This remains open.
\end{conjecture}

\subsection{Infinity}\label{sec:infinity}

The axiom of infinity in ZFC asserts the existence of a
completed infinite set. The standard constructivist
objection is that only potential infinity --- the unbounded
but never completed process --- is legitimate. We argue
that within the discrimination framework, this distinction
is not a valid runtime distinction under the
operationalist axiom (Definition \ref{def:operationalist}),
and therefore collapses.

\begin{definition}[Discriminative infinity]\label{def:infinity}
An \emph{inexhaustible collection} is a collection $S$ such
that for every $\kappa$-state, there exists a predicate whose
articulation (and therefore whose $\kappa$-evolution) produces
a non-trivial residue over $S$. That is: no finite $\kappa$
renders $S$ fully discriminated.
\end{definition}

\begin{proposition}[Indistinguishability of potential
  and actual infinity]\label{prop:infinity}
Let $S$ be an inexhaustible collection. No discriminator
in any $\kappa$-state separates the following two claims:
\begin{enumerate}[label=(\roman*),nosep]
  \item $S$ is a completed infinite totality (actual infinity).
  \item $S$ is an open-ended process that can always yield
    another element upon further $\kappa$-evolution (potential
    infinity).
\end{enumerate}
By the operationalist axiom (Definition
\ref{def:operationalist}), claims that no discriminator
separates are $\kappa$-equivalent and hence identical by
extensionality (Theorem \ref{thm:ext}).
\end{proposition}

\begin{proof}
Suppose, for contradiction, that a discriminator $d$
separates actual from potential infinity. Then $d$ is a
predicate, and its articulation is a $\kappa$-evolution.
After this evolution, $\kappa$ is finite (Proposition
\ref{prop:monotone}: each evolution adds exactly one
dimension). In the new $\kappa$-state, $S$ is still
inexhaustible (by Definition \ref{def:infinity}: no
finite $\kappa$ fully discriminates $S$). The
discriminator $d$ has not resolved the question; it has
merely added one more dimension to a regime that remains
inadequate to exhaust $S$.

No finite sequence of such evolutions converges, because
each produces a $\kappa$-state that is still finite and
$S$ is still inexhaustible. Therefore no discriminator,
and no finite composition of discriminators, can separate
the two claims. They are indistinguishable under every
$\kappa$-state, and hence identical by $\kappa$-extensionality.
\end{proof}

\begin{remark}
This does not deny the usefulness of reasoning about
completed infinities. Classical mathematics routinely and
productively treats the natural numbers, the reals, and
higher cardinalities as completed totalities. The
framework accommodates this by the same mechanism it
accommodates all predication: asserting ``the natural
numbers exist as a completed set'' is an articulable
predicate, and its articulation is a $\kappa$-evolution.
The resulting discriminator is absorbed into $\kappa$ and
is available for subsequent reasoning.

The framework does not restrict \emph{what can be modeled},
only \emph{how models are described}. A Turing machine is
modelable within the framework ($\kappa$ can represent its
state transitions). A Turing machine can serve as a proof
assistant for classical mathematics, including classical
reasoning about completed infinities. Therefore the
framework can reason about anything classical mathematics
can, via the same indirection that all constructive proof
assistants employ.
\end{remark}

\subsection{The operationalist axiom: defense and boundaries}
\label{sec:operationalist-defense}

The operationalist axiom (Definition \ref{def:operationalist})
is the framework's most philosophically load-bearing
commitment. It powers the infinity result (Proposition
\ref{prop:infinity}), infinite choice (Theorem
\ref{thm:choice}(ii)), and the pro-system interpretation
(Corollary \ref{cor:pro-topos}). This subsection
concentrates the defense.

\begin{remark}[Foundational status of indistinguishability]
\label{rem:foundational-status}
Proposition \ref{prop:infinity} is a theorem of the
axiomatized framework: it follows from the operationalist
axiom (Definition \ref{def:operationalist}) applied to
inexhaustible collections. The operationalist axiom plays
the same foundational role that the Axiom of Infinity plays
in ZFC: it is a \emph{choice} about what the framework
admits as meaningful, not a forced mathematical conclusion.
The proof structure is: (1) no discriminator separates the
two claims (proved by contradiction), (2) therefore by the
operationalist axiom (Definition \ref{def:operationalist}), they are identical. Step (2) is an
application of the axiom, not a mathematical derivation.

A Platonist may accept the proof's mechanics entirely ---
including the proof-assistant observation (Remark
\ref{rem:proof-assistant}) that syntactic manipulation of
infinity is finite --- and still maintain that actual and
potential infinity are ontologically distinct. The framework
does not refute Platonism; it \emph{declines to encode it},
on the grounds that ontological distinctions with no
observable structural consequences are outside the
framework's scope (Definition \ref{def:operationalist}).
\end{remark}

\begin{remark}[Epistemic asymmetry and the halting problem]
\label{rem:halting}
The indistinguishability result can be sharpened by analogy
with the halting problem. Consider a system (the ``program'')
reasoning over a collection $S$, and an observer who can see
one $\kappa$-evolution ahead of the system. From the system's
interior, the state space of $S$ is indistinguishable from
infinite at every step: there is always another element
available upon further evolution. From the observer's
exterior, the state space is potentially infinite: unbounded
but evaluated only finitely far at any moment.

The distinction between ``truly infinite'' and ``always
one step ahead'' is exactly the distinction between actual
and potential infinity. Since the system cannot observe its
own bound (this would require solving the halting problem
for its own evolution process), the two are operationally
identical. The operationalist axiom
(Definition \ref{def:operationalist}) applies:
if a distinction is not observable, it is not a valid
runtime distinction.
\end{remark}

\begin{remark}[The hyper-recursive tower]
\label{rem:hyper-recursive}
The classical escape from Remark \ref{rem:halting} is:
``I will step outside the system and observe the
distinction from a meta-level.'' The framework absorbs
this move. The meta-level observer is itself a
$\kappa$-regime (a system reasoning about the first
system's $\kappa$-evolution). It faces the same
structural situation: finite dimensions, inexhaustible
target, perpetual non-arrival of the distinction. So
the critic steps outside \emph{that} --- and arrives in
another $\kappa$-regime, subject to the same constraints.

The triangle identity (Definition \ref{def:triangle})
and free meta-discrimination (\S\ref{sec:meta}) provide
this tower without additional construction: the
$\sigma$-perspective of a $\kappa$-regime is itself a
regime, and the $\sigma$-perspective of \emph{that}
regime is another, and so on. At every level, the
structural invariant holds: finite $\kappa$, inexhaustible
collection, distinction not materialized, and the
\emph{same structural reason} for its non-materialization
as at the previous level.

This is the halting problem applied hyper-recursively.
A system that could resolve actual vs.\ potential
infinity for its tenant would need to solve its own
halting problem, which would require a higher system,
which faces the same obstruction. The arithmetical
hierarchy goes all the way up, and at every level the
invariant holds. The distinction between actual and
potential infinity is not suppressed by fiat; it is
unreachable at every level of a tower that the framework
already contains.

The operationalist axiom (Definition \ref{def:operationalist})'s role is then precise: it says
``unreachable at every level means not a valid
distinction.'' This is a foundational commitment, but
one for which the framework supplies structural
motivation at every level of its own recursive
self-embedding.
\end{remark}

\begin{remark}[The proof assistant concession]
\label{rem:proof-assistant}
A classical mathematician may object that operational
indistinguishability from the interior does not establish
ontological identity: perhaps actual and potential infinity
\emph{are} distinct, and our framework merely cannot see
the distinction.

This objection is refuted by the practice of formalized
mathematics itself. Proof assistants (Lean, Coq, Agda)
run on deterministic, finite-state machines. When a
mathematician verifies Cantor's theorem in Lean, the
computer does not instantiate an uncountably infinite set
in memory. It handles actual infinity \emph{syntactically}:
the Axiom of Infinity is a finite set of rules for symbol
manipulation, and the ``actual infinity'' is a node in a
finite proof tree.

In the discrimination framework, typing the Axiom of
Infinity into a proof assistant is a $\kappa$-evolution:
the articulation of a predicate that adds a discriminator
to a finite regime. The proof assistant then evaluates
theorems about the infinite set by passing them through
this finite $\kappa$-state --- exactly as the framework
prescribes.

By accepting the proof assistant's verdict, the classical
mathematician endorses the premise that every rigorous
truth about actual infinity can be completely captured by
a finite, step-wise process. The only retreat is Platonism:
the claim that while the proof's syntax is finite and
potential, the semantic meaning it points to is actually
infinite. But this is a distinction with no behavioral
consequence: the Platonic infinity produces no observable
structural residue that the syntactic infinity does not
also produce. By the operationalist axiom
(Definition \ref{def:operationalist}),
a distinction with no observable consequence is not a
valid runtime distinction.

The proof assistant is itself a $\kappa$-regime: it has
finite dimensions, it evolves by consuming axioms
(residues from the gap between what is provable and what
is asserted), and it cannot distinguish actual from
potential infinity in its own operation. The framework
does not sidestep the formalization of infinity; it
\emph{describes exactly what proof assistants already do}
when they handle infinity, and names the mechanism.
\end{remark}

\begin{remark}[The framework models the open world, not
  excludes it]
\label{rem:cwa-owa}
The operationalist axiom operates under the Closed World
Assumption (CWA): what is not witnessed by any discriminator
is not a valid distinction. A critic operating under the
Open World Assumption (OWA) may posit unwitnessed structure
beyond the framework's reach.

The framework's response is not exclusion but
\emph{modeling}. The coverage bit $\gamma$
(Section \ref{sec:expanded}) encodes ``not yet
interrogated'' as a positive state in $\GF(2)^3$:
$\gamma = 0$ represents the region that an OWA ontology
would populate with unwitnessed structure. This region is
structurally present, operationally inert until
interrogated, and precisely delineated by the Fano closure
(Theorem \ref{thm:fano}): three independent interrogations
characterize the full shape of what remains opaque.

The recursive tower (Remark \ref{rem:hyper-recursive})
makes this self-similar: the outer system can see the
inner system's $\gamma = 0$ region --- it ``knows what the
inner system doesn't know'' --- but faces its own
$\gamma = 0$ at the next level. The epistemic horizon is
characterizable, invariant across levels, and \emph{that
invariance} is the structural content. Not ``there is
nothing outside'' but ``the outside has this shape from
every inside, and that shape is the same shape at every
level.''

The operationalist axiom is the claim that this
self-similar epistemic horizon is the finest grain at
which the inside/outside distinction can be resolved.
The burden falls on the OWA critic to exhibit a behavioral
residue --- an observable consequence that survives the
collapse --- or accept that the CWA is the appropriate
foundational stance for a framework whose purpose is to
track what can be constructively witnessed.
\end{remark}

\subsection{Choice}\label{sec:choice}

The axiom of choice in ZFC asserts that for every collection
of non-empty sets, there exists a function selecting one
element from each. Its independence from ZF (without Choice)
has been a central topic in set theory since Cohen's forcing
arguments. In the discrimination framework, choice is a
consequence of operational exhaustivity (Theorem
\ref{thm:exhaustive}).

\begin{theorem}[Choice from residue consumption]\label{thm:choice}\footnote{Axiom class \textsf{AO} $\cdot$ depth 7 $\cdot$ \S\ref{app:choice-full}.  \emph{Warrant}: the operationalist axiom collapses the distinction between ``completed selection function'' and ``constructively obtainable for every finite sub-family.'' \emph{Qualifier}: the finite case is class \textsf{A} (constructive). \emph{Rebuttal}: a Platonist may reject the axiom; the finite case stands independently. The operationalist axiom is doing the work that ZFC's axiom of choice does. [O3]}
Let $\mathcal{F} = \{S_i\}_{i \in I}$ be a collection of
non-empty sub-collections of a population under regime
$\kappa$, where the discriminators defining the boundaries
of the $S_i$ have achieved Boolean status
(Definition \ref{def:boolean-dim}) on the relevant
dimensions --- that is, $\chi_d = 1$ for every element
on every dimension that distinguishes the $S_i$ from
each other. Then:
\begin{enumerate}[label=(\roman*),nosep]
  \item For any finite $\mathcal{F}$, there exists a
    constructive $\kappa$-evolution that produces a
    selection function.
  \item For infinite $\mathcal{F}$, no discriminator
    separates ``a completed selection function exists''
    from ``for every finite sub-family, a selection
    function is constructively obtainable.'' By the
    operationalist axiom (Definition
    \ref{def:operationalist}), these claims are
    $\kappa$-equivalent.
\end{enumerate}
Claim (ii) does not prove that choice ``holds'' over
infinite collections in the classical sense. It proves
that the distinction between choice-holding and
choice-not-holding is not a valid runtime distinction
within the framework. A Platonist may accept (i) and
reject (ii) by rejecting the operationalist axiom (Definition \ref{def:operationalist}). (Full analysis: \S\ref{app:choice-full}.)
If the boundaries of the $S_i$ are not Boolean (some
elements have gaps or overlaps on the defining dimensions),
the selection function is not well-defined: the sub-collections
themselves are not sharply delineated, and attempting to
select from them produces a residue rather than a choice.
\end{theorem}

\begin{proof}
\emph{Finite case.}
A selection function over a finite $\mathcal{F} =
\{S_1, \ldots, S_m\}$ requires, for each $S_i$, a
discriminator that isolates a single element of $S_i$
from the rest. Suppose that under $\kappa$, some $S_i$
resists such linearization: no discriminator in $\kappa$
isolates a unique element.

This failure is a discrimination: the wedge product of
the candidate selections over $S_i$ produces a non-zero
residue (the elements that remain indistinguishable). By
Theorem \ref{thm:exhaustive}, this residue fully specifies
its own shape --- it is the blueprint for the discriminator
needed to further separate the candidates. Consuming the
residue via $\kappa$-evolution adds precisely this
discriminator.

Iterate: at each step, either an element is isolated
(selection succeeds for $S_i$) or a non-trivial residue
is produced and consumed, strictly increasing $\dim(\kappa)$
and strictly reducing the number of indistinguishable
candidates. Since each $S_i$ contains finitely many
$\kappa$-equivalence classes (at any finite $\kappa$-state),
and the number of classes strictly decreases with each
evolution, the process terminates for each $S_i$ in
finitely many steps. By Theorem \ref{thm:convergence},
each dimension requires at most $3$ linearly independent
interrogations to achieve full sheaf coverage. Over $m$
finite sub-collections, the total number of evolutions
is finite and bounded.

\emph{Infinite case.}
For infinite $\mathcal{F}$, the iterative process does
not terminate in finitely many steps. However, by
Proposition \ref{prop:infinity} (indistinguishability of
potential and actual infinity), the distinction between
``the selection function exists as a completed object''
and ``for any finite sub-family $\mathcal{F}' \subset
\mathcal{F}$, the selection function over $\mathcal{F}'$
is constructively obtainable'' is not a valid runtime
distinction: no discriminator separates the two claims.

This avoids impredicativity: we do not quantify over
``all discriminators that could be performed'' (which
would be circular). Instead, we observe that the
constructive finite case extends to the infinite case
via the same indistinguishability argument that validates
the axiom of infinity (Proposition \ref{prop:infinity}).
Choice over infinite collections is available to the same
degree that infinite collections themselves are available
(both depend on the operationalist axiom)
--- and as noted in Remark \ref{rem:proof-assistant},
this is exactly the degree to which proof assistants
handle infinite collections when verifying classical
theorems.
\end{proof}

\begin{remark}
A collection resists linearization only under a
\emph{truncation} of $\kappa$. The ``full $\kappa$'' for
a finite sub-collection is not a vague limit notion; it
is the constructive result of exhaustive adjunction
interrogation: $3$ linearly independent interrogations
per dimension (Theorem \ref{thm:convergence}), yielding
a fully glued sheaf (Proposition \ref{prop:presheaf-sheaf}).
For infinite collections, this constructive process
extends by indistinguishability (Proposition
\ref{prop:infinity}).

This is the discrimination framework's account of the
independence of Choice from ZF: in a truncated regime
(analogous to a model of ZF without Choice), selection
functions may not exist because the necessary
discriminators have been projected away. The residue of
the failed selection is the precise structure needed to
recover the missing discriminators. Choice fails only
when you have deliberately erased the information that
would provide it.
\end{remark}

\begin{remark}
This result follows the same pattern as foundation
(Theorem \ref{thm:foundation}), Russell's dissolution
(Theorem \ref{thm:russell}), and infinity (Proposition
\ref{prop:infinity}): what ZFC treats as an axiom, the
discrimination framework derives from the algebraic
structure of $\bigwedge(\GF(2)^n)$, the
residue-consumption dynamics of $\kappa$, and the
regulatory axioms (profile linearity (Definition \ref{def:profile-linear}), the operationalist
axiom).
\end{remark}

\begin{remark}[Choice as a dependent type via Diaconescu]
\label{rem:diaconescu}
In standard topos theory, Diaconescu's theorem establishes
that the Axiom of Choice implies the Law of Excluded Middle
globally. In a topos with a four-valued subobject classifier
$\Omega = \GF(2) \times \GF(2)$, global Choice would force
$\Omega$ to collapse to $\{(1,0), (0,1)\}$, destroying the
four-valued foundation.

In the discrimination framework, this is not a contradiction
but a \emph{typing constraint}. Because Booleanness is a
local dependent type (Proposition \ref{prop:bool-dependent})
rather than a global substrate, Choice cannot be a global
axiom. A selection function exists only over collections
whose defining discriminators have achieved exclusive,
exhaustive partition ($\chi_d = 1$) on the relevant
dimensions. To claim Choice globally would be to claim
Booleanness globally, which the four-valued internal logic
explicitly rejects.

Diaconescu's theorem thus becomes a corroborating feature:
it confirms that the framework's scoping of Choice to
Boolean dimensions is \emph{necessary} for consistency
with the four-valued subobject classifier. Choice, like
classical logic, is earned (constructively verified) per dimension. On dimensions
where Booleanness has not been earned, attempting to
select produces a residue (a gap or overlap in the
sub-collection boundaries), not a selection function. The
residue signals that the preconditions for Choice have
not been met, and $\kappa$-evolution is required before
selection can proceed.
\end{remark}

\begin{example}[Banach--Tarski is non-paradoxical]\footnote{Axiom class \textsf{A} $\cdot$ depth 2 $\cdot$ \S\ref{app:banach-tarski}.  \emph{Warrant}: the construction goes through; the framework does not force the interpretation under which it would be paradoxical. \emph{Qualifier}: this re-describes the standard observation that the pieces are non-measurable; the framework adds the four-cell tracking of which dimensions have earned Booleanness. Finite-$\kappa$: N/A (requires continuous groups). Re-description (standard non-measurability observation, with four-cell tracking).}
\label{ex:banach-tarski}
The Banach--Tarski paradox states that a solid sphere can
be decomposed into finitely many pieces and reassembled
into two spheres identical to the original, using only
rigid motions. This appears to violate conservation of
volume. In the discrimination framework, the construction
goes through --- it is not blocked --- but it is
non-paradoxical, because the framework does not force the
interpretation under which it would be paradoxical.

The classical construction proceeds as follows: the
rotation group $SO(3)$ contains a free subgroup $F_2$ on
two generators. $F_2$ acts on $S^2$ (minus a countable
set of poles) with orbits that partition the sphere. The
axiom of choice selects one representative from each
orbit. The pieces of the decomposition are defined by
which elements of $F_2$ carry the representative to a
given point, classified by the first letter of the
group element's reduced word.

In the discrimination framework, this construction is a
sequence of $\kappa$-evolutions:
\begin{enumerate}[label=(\roman*),nosep]
  \item A discriminator for the free group structure of
    $F_2 \subset SO(3)$: this classifies rotations by
    their reduced-word decomposition.
  \item A discriminator for orbit membership: this
    classifies points of $S^2$ by which $F_2$-orbit
    they belong to.
  \item A choice discriminator selecting one
    representative per orbit (Theorem \ref{thm:choice}):
    this is a $\kappa$-evolution consuming the residue of
    the orbit-selection problem.
  \item Discriminators $d_A, d_B$ assigning each point
    to ``sphere $A$'' or ``sphere $B$'' based on the
    first letter of the group element connecting the
    point to its orbit representative.
\end{enumerate}

The resulting pieces are non-measurable. In the
discrimination framework, this means: the measure
discriminator $d_\mu$ produces $(0,0)$ (gap) on these
pieces. The gap arises because $d_\mu$ is defined by
geometric predicates (countable unions of intervals,
limits of simple functions) while the pieces are defined
by algebraic predicates (group-word classification,
orbit-representative selection). These two families of
discriminators operate on \emph{algebraically independent
dimensions of $\kappa$}: the wedge product
$d_\mu \wedge d_A$ is non-zero (they are linearly
independent), confirming that they discriminate
independently and do not constrain each other.

Volume ``doubling'' is the artefact of assuming that
$d_\mu$ achieves Boolean status
(Definition \ref{def:boolean-dim}) over all subsets,
including those defined by algebraically independent
discriminators. The framework reports: $d_\mu$ and
$d_A$ do not jointly achieve Booleanness on the
decomposed pieces. The $(0,0)$ gap on $d_\mu$ is the
framework correctly recording that measure has no
verdict on sets constructed by group-theoretic choice
--- not because measure ``fails,'' but because the
geometric and algebraic dimensions of $\kappa$ are
independent, and Booleanness across independent
dimensions must be earned, not assumed.
\end{example}

